\documentclass[12pt,a4paper]{article}

\usepackage{fontawesome5}
\providecommand{\faPencilAlt}{\faIcon{pencil-alt}}

\usepackage{styles/couleurs}
\usepackage{styles/mise_en_page}
\usepackage{styles/boites}
\usepackage{styles/maths}
\usepackage{styles/figures}

\usepackage{amsmath}
\usepackage{enumitem}
\usepackage{tcolorbox}

\geometry{
  a4paper,
  top=1.2cm,
  bottom=1.2cm,
  left=1.4cm,
  right=1.4cm
}
\pagestyle{empty}

\begin{document}

\begin{center}
  {\sffamily\LARGE\bfseries\color{TitreC}
  Point rédaction -- Calculer une moyenne}
\end{center}

\vspace{-0.25em}

\remarque{Une moyenne est un \textbf{quotient} :
on additionne toutes les valeurs, puis on divise par le nombre de valeurs.

Dans une chaîne de calculs, le signe \(=\) relie toujours des expressions
qui ont la \textbf{même valeur}.}

\begin{tcolorbox}[
  mathsbox,
  colback=SavoirBg,
  colframe=SavoirC,
  boxrule=0.7pt,
  arc=2mm,
  left=3mm,
  right=3mm,
  top=2mm,
  bottom=2mm
]
{\sffamily\bfseries\color{SavoirC}
Méthode 1 -- En une seule expression}

\smallskip

Pour les valeurs
\[
7 \ ;\ 12 \ ;\ 10 \ ;\ 5 \ ;\ 9 \ ;\ 8,
\]
on peut écrire avec une \textbf{barre de fraction} :

\[
M=
\frac{7+12+10+5+9+8}{6}
=
\frac{51}{6}
=
8{,}5.
\]

Ou avec le symbole \(\div\) :

\[
M=
(7+12+10+5+9+8)\div 6
=
51\div 6
=
8{,}5.
\]

\textbf{Attention :} avec \(\div\), les parenthèses permettent de montrer
que l'on divise \textbf{toute la somme} par \(6\).

\medskip
\textbf{À privilégier :} cette écriture permet de s'entraîner à lire
et à manipuler une expression numérique complète.
\end{tcolorbox}

\begin{tcolorbox}[
  mathsbox,
  colback=NeutreBg,
  colframe=NeutreC!55,
  boxrule=0.6pt,
  arc=2mm,
  left=3mm,
  right=3mm,
  top=2mm,
  bottom=2mm
]
{\sffamily\bfseries\color{NeutreC}
Méthode 2 -- En deux étapes}

\smallskip

Cette rédaction est \textbf{tout aussi correcte} :

\[
7+12+10+5+9+8=51
\]

Il y a \(6\) valeurs.

\[
M=51\div 6=8{,}5.
\]

\end{tcolorbox}

\begin{tcolorbox}[
  mathsbox,
  colback=ActionBg,
  colframe=ActionC,
  boxrule=0.7pt,
  arc=2mm,
  left=3mm,
  right=3mm,
  top=2mm,
  bottom=2mm
]
{\sffamily\bfseries\color{ActionC}
Et avec un tableau d'effectifs ?}

\smallskip

Même principe : chaque valeur est multipliée par son effectif.

\[
M=
\frac{35\times7+36\times9+37\times5+38\times2+39\times7}
     {7+9+5+2+7}
=
\frac{1103}{30}
\approx
36{,}8.
\]

\end{tcolorbox}

\end{document}
